\section{Ambiguous Specifications and Lifted Monotonicity}
\label{sec:ambiguous-specifications}

A functional specification fixes one output at each input.  We now allow a
fiber to contain several outputs and ask when they can be preserved together
as input grows.  Two examples introduce the issue; the formal development then
extends the specification, monotonicity condition, and transition system from
Section~\ref{sec:functional-calm}.

\subsection{The framework on two ambiguous examples}
\label{sec:ambiguity-story}

We begin, as in Section~\ref{sec:election-story}, with an informal,
example-driven treatment.  Terms introduced in italics are defined formally
below, and nothing here is needed for the proofs.  Section~\ref{sec:election-story}
exercised the shared transition system; these examples exercise the additions
for ambiguous specifications: directedness, output-choice commitments, output
retirement, and Lifted Monotonicity.

\begin{figure}[t]
  \centering
  \small
  \setlength{\tabcolsep}{5pt}
  \begin{tabular}[t]{@{}l l@{}}
    \multicolumn{2}{@{}l}{(f) Transaction ordering, after the commitment to $[A,B]$}\\
    \toprule
    Input & Output \\ \midrule
    $\{A,B\}$ & \sh{$[A,B]$} \\
    \retired{$\{A,B\}$} & \retired{$[B,A]$} \\ \strikerow
    \bottomrule
    \multicolumn{2}{@{}l}{\footnotesize no remaining row refines the removed one:}\\
    \multicolumn{2}{@{}l}{\footnotesize a commitment}
  \end{tabular}
  \hfill
  \begin{tabular}[t]{@{}l l@{}}
      \multicolumn{2}{@{}l}{(g) Traceroute, after outputting $\{s{\to}\texttt{s1}\}$}\\
    \toprule
    Input & Output \\ \midrule
    $i$ & \sh{$\{s{\to}\texttt{s1}\}$} \\
    \retired{$i$} & \retired{$\{s{\to}a\}$} \\
    \retired{$i$} & \retired{$\{s{\to}a,\ s{\to}a{\to}\texttt{s2}\}$} \\
    $i$ & \sh{$\{s{\to}\texttt{s1},\ s{\to}a\}$} \\
    $i$ & \sh{$\{s{\to}\texttt{s1},\ s{\to}a,\ s{\to}a{\to}\texttt{s2}\}$} \\
    \bottomrule
    \multicolumn{2}{@{}l}{\footnotesize each removed row refines to a remaining one:}\\
    \multicolumn{2}{@{}l}{\footnotesize exposable}
  \end{tabular}
  \caption{Two specifications with several rows at one current input, each
  shown after one output action.  In (f) the input is a pair of
  concurrent transactions and an output is a serial order of them.  In (g) an
  output is a prefix-closed set of routes from source $s$, which reaches server
  \texttt{s1} directly and server \texttt{s2} through hop $a$.  Shaded cells
  are the upper fiber that remains.  Legend as in
  Figure~\ref{fig:election-frames}.}
  \label{fig:ambiguity-frames}
\end{figure}

\paragraph{Output actions under choice.}
As in the functional case, outputting $o$ promises that the emitted result will
be refined by the final result.  The single atomic action
$\mathsf{output}(o)$ applies any required commitments, retires dominated rows,
and emits $o$.  With several rows at one input, retirement may remove a row
without losing an alternative: a retained row at the same input may refine both
$o$ and the removed output.  Output $o$ is \emph{exposable} when its action needs
no commitment.

\paragraph{Ambiguity.}
Figure~\ref{fig:ambiguity-frames}(f) allows either serial order of concurrent
transactions $A$ and $B$.  No output refines both orders, so outputting one
requires an \emph{output-choice commitment} that removes the other.  In the
traceroute of panel (g), by contrast, outputs are prefix-closed route sets and
their unions are within spec.  Outputting $\{s{\to}\texttt{s1}\}$ requires no
commitment; the action retires rows whose outputs do not refine it, but a
retained union refines each retired output.  Nothing is lost, so the output is
exposable.

The distinction is \emph{directedness}: every finite set of outputs in a fiber
must have a common refinement in that fiber.  A directed fiber allows any output
to be output without losing an alternative; a nondirected fiber may require
an output-choice commitment.  Lifted Monotonicity combines this requirement with the
requirement from Functional CALM that more input not invalidate an output:
every finite set of outputs within spec at $i$ must have a common refinement
within spec at every $j\succeq_I i$.

\paragraph{The commitment basis.}
The specification fixes which rows conflict with an output, while the system's
commitment basis determines which required seal and output-choice commitments
it can perform and how it bundles them.  In Figure~\ref{fig:cc-commitments},
two-phase locking commits only $A$ before $B$, whereas multiversion timestamp
ordering commits $B$ to first place.  Appendix~\ref{app:cc-bases} gives both
bases, including their abort commitments, and
Section~\ref{sec:design-boundary} returns to this design choice.

\begin{figure}[t]
  \centering
  \small
  \setlength{\tabcolsep}{5pt}
  \begin{tabular}[t]{@{}l l@{}}
    \multicolumn{2}{@{}l}{(h) 2PL: $A$ locks ahead of $B$}\\
    \toprule
    Input & Output \\ \midrule
    $\{A,B,C\}$ & \sh{$[A,B,C]$} \\
    $\{A,B,C\}$ & \sh{$[A,C,B]$} \\
    $\{A,B,C\}$ & \sh{$[C,A,B]$} \\
    \retired{$\{A,B,C\}$} & \retired{$[B,A,C]$} \\ \strikerow
    \retired{$\{A,B,C\}$} & \retired{$[B,C,A]$} \\ \strikerow
    \retired{$\{A,B,C\}$} & \retired{$[C,B,A]$} \\ \strikerow
    \bottomrule
  \end{tabular}
  \hfill
  \begin{tabular}[t]{@{}l l@{}}
    \multicolumn{2}{@{}l}{(i) MVTO: $B$ receives the first timestamp}\\
    \toprule
    Input & Output \\ \midrule
    $\{A,B,C\}$ & \sh{$[B,A,C]$} \\
    $\{A,B,C\}$ & \sh{$[B,C,A]$} \\
    \retired{$\{A,B,C\}$} & \retired{$[A,B,C]$} \\ \strikerow
    \retired{$\{A,B,C\}$} & \retired{$[A,C,B]$} \\ \strikerow
    \retired{$\{A,B,C\}$} & \retired{$[C,A,B]$} \\ \strikerow
    \retired{$\{A,B,C\}$} & \retired{$[C,B,A]$} \\ \strikerow
    \bottomrule
  \end{tabular}
  \caption{Serialization orders for three concurrent transactions, after one
  commitment under two protocols.  In (h) two-phase locking has ordered $A$
  before $B$; in (i) multiversion timestamp ordering has placed $B$ first.
  Legend as in Figure~\ref{fig:election-frames}.}
  \label{fig:cc-commitments}
\end{figure}

These examples isolate the new specification-level issue: several outputs may
remain possible at one input, and a system may have to choose among them.  We
now formalize when those alternatives can be preserved together.

\subsection{Relational specifications and Lifted Monotonicity}
\label{sec:relational-specifications}

The examples pose the specification-level question: when can every finite set
of possible outputs survive together as input grows?  We first define the two
ways this can fail, then combine them in Lifted Monotonicity.

\begin{definition}[Ordered relational specification]
\label{def:relational-specification}
An \emph{ordered relational specification} is
        \begin{equation}
  \Spec=(I,\InputOrd,O,\OutputOrd,S),
  \qquad S\subseteq I\times O,
  \label{eq:relational-specification}
        \end{equation}
where $(I,\InputOrd)$ and $(O,\OutputOrd)$ are partially ordered sets.  A row
$(i,o)\in S$ says that output $o$ is within spec at input $i$.  The
\emph{fiber} of $S$, or of any $R\subseteq S$, at $i$ is the set of outputs
paired with $i$,
        \begin{equation}
  R(i)=\{o\in O:(i,o)\in R\},
  \label{eq:fiber}
        \end{equation}
and may now contain several outputs; the upper fiber of
Section~\ref{sec:canonical-transition-system} is the union of the fibers at $i$ and its
valid refinements.  The original specification has a nonempty fiber at every
input under discussion.  A functional specification is the special case
$\fiber{i}=\{f(i)\}$.
\end{definition}

The specification-state and input-event definitions of
Section~\ref{sec:determination} apply unchanged to relational tables.  Output
actions and runs are extended after the Lifted Monotonicity criterion is defined
below.

The output-refinement order lifts from individual outputs to nonempty sets of
outputs.  For nonempty $A,B\subseteq O$, define
        \begin{equation}
  A\LiftedOutputOrd B
  \quad\Longleftrightarrow\quad
  \forall a\in A\;\exists b\in B:\;a\OutputOrd b.
  \label{eq:lifted-output-order}
        \end{equation}
This relation is often called the Hoare preorder.  It compares sets
pointwise: $A\LiftedOutputOrd B$ means that every output in $A$ is
refined by some output in $B$.

\begin{definition}[Relational Monotonicity]
\label{def:relational-monotonicity}
An ordered relational specification satisfies \emph{Relational Monotonicity}
when the mapping $i\mapsto\fiber{i}$ is monotone
under the pointwise lift:
        \begin{equation}
  i\InputOrd j
  \quad\Longrightarrow\quad
  \fiber{i}\LiftedOutputOrd\fiber{j}.
  \label{eq:relational-monotonicity}
        \end{equation}
\end{definition}

For singleton fibers this is functional monotonicity,
Equation~\ref{eq:functional-monotonicity}.  In general it says that each output within spec before more input arrives has
some refinement afterward.  It does not say whether several outputs within
spec at one input can be preserved
together; Directed Outputs addresses that question.

\begin{definition}[Compatibility]
\label{def:compatibility}
Let $F\subseteq O$ be a set of outputs.  Outputs $o_1,\ldots,o_k$ are
\emph{compatible in $F$} if some $u\in F$ refines all of them,
$o_r\OutputOrd u$ for every $r$, and \emph{incompatible in $F$} otherwise.
They are \emph{compatible at input $i$} when they are compatible in the fiber
$\fiber{i}$.
\end{definition}

\begin{definition}[Directed fiber; Directed Outputs]
\label{def:fiber-directedness}
A fiber $F\subseteq O$ is \emph{directed} when every nonempty finite
$X\subseteq F$ is compatible in $F$.  An ordered relational specification has
\emph{Directed Outputs} when every fiber $\fiber{i}$ is directed.
\end{definition}

A singleton fiber is directed, and closure under a join is sufficient but not
necessary.  Directedness requires some upper bound within the fiber for each finite
set, not a least one or a fixed merge operator, so it is weaker than a
join-semilattice.  Compatibility likewise does not require comparability: two
incomparable outputs may have a third output in the fiber above both.  Thus the
traceroute fiber is directed by union, whereas the two transaction orders in
Figure~\ref{fig:ambiguity-frames}(f) are incompatible.

\paragraph{Lifted Monotonicity.}
\label{sec:complete-calm-criterion}

Relational Monotonicity preserves each output across input refinement;
Directed Outputs preserves finite choices within a fiber.  The joint lift
expresses both.  For nonempty $A,B\subseteq O$, define
        \begin{equation}
  A\JointOutputOrd B
  \;\Longleftrightarrow\;
  \text{every nonempty finite }X\subseteq A\text{ is compatible in }B,
  \label{eq:joint-output-order}
        \end{equation}
the \emph{joint lift} of $\OutputOrd$.  Unlike the Hoare preorder, it
preserves finite sets of outputs together rather than each output separately.
It is transitive, and $A\JointOutputOrd A$ holds exactly when $A$ is directed.

\begin{definition}[Lifted Monotonicity; monotone specification]
\label{def:lifted-monotonicity}
An ordered relational specification satisfies \emph{Lifted Monotonicity}, and
is called \emph{monotone}, when the mapping $i\mapsto\fiber{i}$ is monotone
under the joint lift:
        \begin{equation}
  i\InputOrd j
  \quad\Longrightarrow\quad
  \fiber{i}\JointOutputOrd\fiber{j}.
  \label{eq:lifted-monotonicity}
        \end{equation}
\end{definition}

The next proposition supplies two forms used later: form~(2) separates the
across-input and within-fiber requirements, and form~(3) matches confluence by
allowing outputs from several inputs.

\begin{proposition}[Equivalent forms of Lifted Monotonicity]
\label{prop:lifted-monotonicity-decomposition}
For an ordered relational specification, the following are equivalent:
\begin{enumerate}
  \item Lifted Monotonicity: $i\InputOrd j$ implies
        $\fiber{i}\JointOutputOrd\fiber{j}$;
  \item Relational Monotonicity and Directed Outputs;
  \item for every finite set of rows $(i_1,o_1),\ldots,(i_k,o_k)$ of $S$, at
        inputs that need not be comparable, the outputs $o_1,\ldots,o_k$ are
        compatible at every common refinement of $i_1,\ldots,i_k$.
\end{enumerate}
\end{proposition}

\begin{proof}
(1) implies (2).  Taking $j=i$, $\fiber{i}\JointOutputOrd\fiber{i}$ says that
$\fiber{i}$ is directed.  Taking $X=\{o\}$ for $o\in\fiber{i}$ and
$i\InputOrd j$ gives $u\in\fiber{j}$ with $o\OutputOrd u$, which is Relational
Monotonicity.

(2) implies (3).  Let rows $(i_r,o_r)$ have common refinement $j$.  Since
$i_r\InputOrd j$, Relational Monotonicity gives, for each $r$, a
$u_r\in\fiber{j}$ with $o_r\OutputOrd u_r$.  The $u_r$ now all lie in the one
fiber $\fiber{j}$, which is directed, so some $u\in\fiber{j}$ refines every
$u_r$, hence every $o_r$.

(3) implies (1).  Let $i\InputOrd j$ and $X\subseteq\fiber{i}$ be nonempty and
finite.  The rows $\{(i,x):x\in X\}$ have common refinement $j$, so some
$u\in\fiber{j}$ refines all of $X$; that is $\fiber{i}\JointOutputOrd\fiber{j}$.
\end{proof}

For singleton fibers this is functional monotonicity.  Form~(3) quantifies over
all finite sizes because pairwise compatibility is insufficient
(Appendix~\ref{app:additional-applications}), and over incomparable inputs
because different runs meet only at a common refinement.

\subsection{Output actions under choice}
\label{sec:safe-exposure}

The functional case has one row per input, so an output action either removes
nothing or seals a valid input.  With several rows at one input, it can also
retire a dominated output while retaining another row that preserves its
claim.  We now extend output actions to this case.

For any remainder table, write
$R(k)=\{u\in O:(k,u)\in R\}$.  For $o\in R(i)$, let
\begin{equation}
  R_o=R\mathbin{\upharpoonright_O}o
      =\{(k,u)\in R:o\OutputOrd u\}
  \label{eq:relational-exposure-target}
\end{equation}
be the target of the canonical output action.  The rows in
$R\setminus R_o$ fall into three classes:
\begin{align}
  K_o&=\{k:R(k)\neq\varnothing\text{ and }R_o(k)=\varnothing\},
      \label{eq:seal-inputs}\\
  D_o&=\{(k,x)\in R\setminus R_o:
          \exists u\in R_o(k).\ x\OutputOrd u\},
      \label{eq:output-retirements}\\
  U_o&=(R\setminus R_o)
       \setminus\bigl(D_o\cup\{(k,x):k\in K_o\}\bigr).
      \label{eq:choice-removals}
\end{align}
Rows at inputs in $K_o$ require seal commitments.  Rows in $D_o$ are dominated
by rows retained at the same input and are retired by the output action.  Rows
in $U_o$ are undominated conflicts with $o$ and require an
\emph{output-choice commitment}.  Unlike a seal commitment, an output-choice
commitment may remove some outputs while leaving their input valid.

The canonical output action first seals the inputs in $K_o$, then removes
$U_o$.  Concretely, let
\[
  C_o^{\mathrm{seal}}=\{(k,x)\in R:k\in K_o\},
  \qquad R_o^{\mathrm{seal}}=R\setminus C_o^{\mathrm{seal}}.
\]
The canonical seal commitment is
\begin{equation}
  s_{(i,R),o}(i,R)=(i,R_o^{\mathrm{seal}}),
  \label{eq:canonical-seal-commitment-relational}
\end{equation}
when $K_o\neq\varnothing$.  The canonical output-choice commitment is
\begin{equation}
  c_{(i,R),o}\bigl(i,R_o^{\mathrm{seal}}\bigr)
    =\bigl(i,R_o^{\mathrm{seal}}\setminus U_o\bigr),
  \label{eq:canonical-output-choice-commitment}
\end{equation}
when $U_o\neq\varnothing$.  The canonical commitment basis contains these
nonidentity commitments for every state and permitted output.  After applying
them internally, $\mathsf{output}(o)$ emits $o$ and retires $D_o$, so the atomic
action ends at $(i,R_o)$.

\begin{definition}[Output retirement, under output choice]
\label{def:output-retirement}
The \emph{output retirements} for outputting $o$ at $(i,R)$ are the rows in
$D_o$: rows excluded by $o$ but covered by a retained row at the same input.
\end{definition}

An output action under another basis is \emph{legal} when its commitments retain
$(i,o)$ and, after those commitments, every remaining row excluded by $o$ is
dominated at the same input by a retained row that refines $o$; the action then
retires those dominated rows.  The canonical basis is maximally preserving:

\begin{lemma}[Canonical commitment normal form]
\label{lem:canonical-commitment-normal-form}
The canonical basis supports an output action for every permitted output.  Every
legal output action for $o$ removes the canonical seal and output-choice rows
for $o$.  The canonical action $\mathsf{output}(o)$ therefore emits $o$ and
retains every row not required to honor that output.  Any additional row removal
only narrows the resulting remainder table.
\end{lemma}

\begin{proof}
Every row at an input in $K_o$ conflicts with $o$, so that input must be sealed.
At every remaining input, each row in $U_o$ conflicts with $o$ and is not
covered by a row in $R_o$, so it cannot be retired and must be removed by a
commitment.  The rows in $D_o$ are covered and retired.  These three sets
partition $R\setminus R_o$, so the canonical action ends at exactly $R_o$.  Any
action that removes more rows ends at a subset of $R_o$ and only excludes
additional behavior.
\end{proof}

The canonical seal and output-choice commitments above extend the commitment
basis in Definition~\ref{def:transition-system}.  The only specification state
is $\sigma=(i,R)$; there is no pending-output marker.  A relational run uses one
atomic event
\[
  (i,R)\xrightarrow{\mathsf{output}(o)}(i,R'),
\]
which atomically applies its commitments, retires the remaining covered rows,
and emits $o$.  Under the canonical basis $R'=R_o$.  The transition-system and
run definitions,
Lemma~\ref{lem:run-residual}, Consistency, and
Proposition~\ref{prop:confluence} therefore extend unchanged to relational
specifications.

\begin{definition}[Exposable, under output choice]
\label{def:exposable-relational}
An output $o\in R(i)$ is \emph{exposable} at $(i,R)$ when a legal
$\mathsf{output}(o)$ action requires no commitment; it may still retire
dominated rows.
\end{definition}

For singleton fibers $D_o$ and $U_o$ are empty, so this reduces to
Definition~\ref{def:exposable}.

\begin{proposition}[Exposability]
\label{prop:exposability}
An output $o\in R(i)$ is exposable at $(i,R)$ if and only if, for every valid
$j\succeq_I i$ and $x\in R(j)$, some $u\in R(j)$ refines both:
\begin{equation}
  o\OutputOrd u
  \qquad\text{and}\qquad
  x\OutputOrd u.
  \label{eq:safe-exposure}
\end{equation}
When $o$ is exposable, $\mathsf{output}(o)$ uses no commitment and retires only
dominated rows.
\end{proposition}

\begin{proof}
At a valid $j$, the canonical output action retains
$R_o(j)=\{u\in R(j):o\OutputOrd u\}$.  A removed $x$ can be retired exactly when
some $u\in R_o(j)$ refines it, which is the stated test; the same witness makes
$R_o(j)$ nonempty.  Hence a commitment-free output action is legal exactly under
that test.
\end{proof}

This test marks the coordination boundary: if it fails, every legal output
action uses a commitment to remove conflicting rows.  Section~\ref{sec:complete-calm}
combines this boundary with the distributed execution model.
